\documentclass{article}

% Margins and spacing
\usepackage[a4paper, margin=.7in]{geometry}
% Math packages
\usepackage{amsmath}
\usepackage{amssymb}
% Placement of figures, tables, algorithms, etc.
\usepackage{float}
% Algorithms
\usepackage{algorithm}
\usepackage{algpseudocode}

% No section numbering
\setcounter{secnumdepth}{0}
% No paragraph indentation
\setlength{\parindent}{0pt}

% Define custom commands
%% Operators
\DeclareMathOperator*{\argmin}{arg\,min}
\DeclareMathOperator*{\argmax}{arg\,max}
%% Manually indented state in algorithmic environment
\newcommand{\IndentState}{\State \hspace{1em}}

\title{CSCI-570 Homework 4 Solutions}
\author{
  Naman Chhibbar \\
  USC ID: 7709-8548-02
}
\date{}

% Document body
\begin{document}
\maketitle

\section{Problem 1}
  Let $\textsc{SubsetSum}(A, m, n)$ be the required algorithm for an array $A$ of size $m$ and target sum $n$.
  The algorithm returns \texttt{True} if there exists a subset of $A$ that sums to $n$, and \texttt{False} otherwise.
  Note that this problem can be easily broken down into subproblems as:
  \[
    \textsc{SubsetSum}(A, m, n) =
    \underbrace{\textsc{SubsetSum}(A[1:m-1], m-1, n)}_{\text{exclude } A[0]}
    \lor \underbrace{\textsc{SubsetSum}(A[1:m-1], m-1, n - A[0])}_{\text{include } A[0]}
    \tag{1}
  \]
  where $A[i:j]$ is the subarray of $A$ from index $i$ to $j$ (both inclusive) while following 0-based indexing and $\lor$ denotes the logical or operation (we assume $A$ to be of length more than one for generalization; the base cases are covered later).
  Let $M$ be the memory matrix of size $(m+1) \times (n+1)$ defined as $M[i][j] = \textsc{SubsetSum}(A[0:i], i, j)$.
  From (1), we have:
  \[ M[i][j] = M[i-1][j] \lor M[i-1][j - A[i-1]] \]
  Note that the $i$th row of $M$ only depends on the $(i-1)$th row.
  Thus, we can optimize the space complexity of the algorithm by storing only two rows of $M$ at any time.
  The algorithm becomes:
  \begin{algorithm}
    \caption{Subset Sum}
    \begin{algorithmic}[1]
      \State \textbf{Input:} Array $A$ of $m$ positive integers, target sum $n > 0$
      \State \textbf{Output:} \texttt{True} if there exists a subset of $A$ that sums to $n$, \texttt {False} otherwise
      \State Initialize $M$ of size $2 \times (n+1)$ with $M[0][0] = \texttt{True}$ and $M[0][j] = \texttt{False}$ for $j > 0$
      \For{$i = 0 \to (m-1)$}
        \For{$j = 0 \to n$}
          \If{$j < A[i]$}
            \State $M[1][j] \gets M[0][j]$
          \Else
            \State $M[1][j] \gets M[0][j] \lor M[0][j - A[i]]$
          \EndIf
        \EndFor
      \EndFor
      \State \Return $M[1][n]$
    \end{algorithmic}
  \end{algorithm}
  
  We can see that we do a constant amount of work in each iteration of the loops.
  Thus, the time complexity of the algorithm is $O(mn)$ and the space complexity is $O(n)$.
  The time complexity is polynomial in $m$ but, since $n$ is an integer, the time and space complexity is pseudo-polynomial in $n$.

\section{Problem 2}
  Note that the robot can reach the lower right corner by taking exactly $m$ down steps and $n$ right steps in any order ($m$ being the number of rows arranged horizontally and $n$ being the number of columns).
  Hence, the number of possible ways is the number of permutations of:
  \[ \underbrace{\downarrow \downarrow \cdots \downarrow}_{m} \underbrace{\rightarrow \rightarrow \cdots \rightarrow}_{n} \]
  which is $\frac{(m+n)!}{m!n!}$.
  To reduce the number of computations, we can compute it as $\frac{(m+1)(m+2) \cdots (m+n)}{n!}$ if $m \ge n$ or $\frac{(n+1)(n+2) \cdots (m+n)}{m!}$ otherwise.
  The algorithm becomes:
  \begin{algorithm}[H]
    \caption{Number of Paths}
    \begin{algorithmic}[1]
      \State \textbf{Input:} Integers $m, n$
      \State \textbf{Output:} Number of unique paths from top-left to bottom-right of an $m \times n$ grid
      \State $numerator = 1$
      \State $denominator = 1$
      \If{$m \ge n$}
        \For{$i = 1 \to n$}
          \State $numerator \gets numerator \times (m + i)$
          \State $denominator \gets denominator \times i$
        \EndFor
      \Else
        \For{$i = 1 \to m$}
          \State $numerator \gets numerator \times (n + i)$
          \State $denominator \gets denominator \times i$
        \EndFor
      \EndIf
      \State \Return $numerator / denominator$
    \end{algorithmic}
  \end{algorithm}

  Since we do a constant amount of work in each iteration of the loop, the time complexity of the algorithm is $O(\min(m, n))$.

\section{Problem 3}
  Let \textsc{LongestPalindrome}($S, n$) be the required algorithm that returns the length of the longest palindromic subsequence in the string $S$ of length $n$.
  Note that the problem can be broken down as:
  \[
    \textsc{LP}(S, n) =
    \begin{cases}
      \textsc{LP}(S[1:n-2], n-2) + 2 & S[0] = S[-1] \\
      \max\{
        \underbrace{\textsc{LP}(S[0:n-2], n-1)}_{\text{choose } S[0]},
        \underbrace{\textsc{LP}(S[1:n-1], n-1)}_{\text{choose } S[n - 1]}
      \}
        & \text{otherwise}
    \end{cases}
    \tag{1}
  \]
  where $S[i:j]$ is the substring of $S$ from index $i$ to $j$ (both inclusive) while following 0-based indexing and \textsc{LP} denotes the algorithm \textsc{LongestPalindrome} (assuming $S$ is of length more than 2 for generalization; bases cases are covered later).
  Let $M$ be the memory matrix of size $n \times n$ defined cleverly as $M[i][j] = \textsc{LP}(S[j:i+j-1], i)$.
  Let's generalize (1) for indices $i$ to $j$:
  \[
    \textsc{LP}(S[i:j], j-i+1) =
    \begin{cases}
      0 & i > j \\
      1 & i = j \\
      \textsc{LP}(S[i+1:j-1], j-i-1) + 2 & S[i] = S[j] \\
      \max\{
        \underbrace{\textsc{LP}(S[i:j-1], j-i)}_{\text{choose } S[i]},
        \underbrace{\textsc{LP}(S[i+1:j], j-i)}_{\text{choose } S[j]}
      \}
        & \text{otherwise}
    \end{cases}
  \]
  Substituting $i \to j$ and $j \to i + j - 1$, we get:
  \[
    \textsc{LP}(S[j:i+j-1], i) =
    \begin{cases}
      0 & i < 1 \Rightarrow i = 0 \\
      1 & i = 1 \\
      \textsc{LP}(S[j+1:i+j-2], i-2) + 2 & S[j] = S[i+j-1] \\
      \max\{\textsc{LP}(S[j:i+j-2], i-1), \textsc{LP}(S[j+1:i+j-1], i-1)\} & \text{otherwise}
    \end{cases}
    \tag{2}
  \]
  From (2), we can define $M$ as:
  \[
    M[i][j] =
    \begin{cases}
      0 & i = 0 \\
      1 & i = 1 \\
      M[i-2][j+1] + 2 & S[j] = S[i+j-1] \\
      \max\{M[i-1][j], M[i-1][j+1]\} & \text{otherwise}
    \end{cases}
  \]
  Due to the way we defined $M$, we can see that again the $i$th row of $M$ only depends on the rows $(i-1)$ and $(i-2)$.
  Hence, we can optimize the space complexity by storing only three rows of $M$ at a time.
  Also, the required answer is $M[n][0] = \textsc{LongestPalindrome}(S[0:n-1], n)$.
  Note that, since $M[i][j]$ depends on values in column $j$ and $j+1$, we need to iterate $j$ from $n-i$ to $0$.
  Also note that the values of $M$ are not defined for $j > n - i$ (since $i + j - 1 \le n - 1 \Rightarrow j \le n - i$ must hold).
  The algorithm becomes:
  \begin{algorithm}
    \caption{Longest Palindromic Subsequence}
    \begin{algorithmic}[1]
      \State \textbf{Input:} String $S$ of length $n$
      \State \textbf{Output:} Length of the longest palindromic subsequence in $S$
      \State Initialize $M$ of size $3 \times n$ with $M[0][j] = 0$ and $M[1][j] = 1$ for all $j$, $M[1][n] = 0$
      \For{$i = 2 \to n$}
        \For{$j = n-i \to 0$}
          \If{$S[j] = S[i+j-1]$}
            \State $M[2][j] \gets M[0][j+1] + 2$
          \Else
            \State $M[2][j] \gets \max\{M[1][j], M[1][j+1]\}$
          \EndIf
        \EndFor
        \For{$k = 0 \to n$}
        \State $M[0][k] \gets M[1][k]$
        \State $M[1][k] \gets M[2][k]$
        \EndFor
      \EndFor
      \State \Return $M[1][0]$
    \end{algorithmic}
  \end{algorithm}

  Since we have two nested loops doing a constant amount of work in each iteration, the time complexity of the algorithm is $O(n^2)$.
  And the space complexity is $O(n)$ since we are storing only three rows of $M$ at any time.

\section{Problem 4}
  Consider three lines $L_i, L_j, L_k$ with the parameters $(a_i, b_i)$, $(a_j, b_j)$, and $(a_k, b_k)$ respectively.
  Line $L_i$ is visible if $\exists x \in \mathbb{R}$ such that:
  \begin{align*}
    a_i x + b_i &> a_j x + b_j \\
    \Rightarrow (a_i - a_j) x &> b_j - b_i \\
  \end{align*}
  Similarly, we also have $(a_i - a_k) x > b_k - b_i$.
  Assuming $a_j < a_i < a_k$, we have:
  \begin{align*}
    x > \frac{b_j - b_i}{a_i - a_j} \quad &\text{and} \quad x < \frac{b_k - b_i}{a_i - a_k} \\
    \Rightarrow \frac{b_j - b_i}{a_i - a_j} &< x < \frac{b_k - b_i}{a_i - a_k}
  \end{align*}
  Hence, $L_i$ is visible if $\frac{b_j - b_i}{a_i - a_j} < \frac{b_k - b_i}{a_i - a_k}$.
  We can use this as a check to see if a line is visible, given two other lines having slopes like $a_j, a_k$. \\
  Also, note that if we have parallel lines $L_1, \ldots, L_n$, only $L_i$ will be visible where $i = \argmax_{j \in \{1, \ldots, n\}} b_j$. \\
  The required algorithm is:
  \begin{algorithm}[H]
    \caption{Visible Lines}
    \begin{algorithmic}[1]
      \State \textbf{Input:} Array $L$ of $n$ lines where $L[i] = (a_i, b_i)$
      \State \textbf{Output:} Array of visible lines in $L$
      \State Sort $L$ in increasing order of $a_i$
      \State Initialize new empty array $L'$
      \Comment To remove lines with same slopes
      \State $L'[0] \gets L[0]$
      \State $n' = 1$
      \Comment Length of $L'$
      \State $slope = L[0][0]$
      \For{$i = 1 \to (n-1)$}
        \If{$L[i][0] \ne slope$}
          \State $L'[n'] \gets L[i]$
          \State $slope \gets L[i][0]$
          \State $n' \gets n' + 1$
        \ElsIf{$L[i][1] > L'[n'-1][1]$}
          \State $L'[n'-1] \gets L[i]$
        \EndIf
      \EndFor
      \State Initialize empty array $V$
      \Comment To store visible lines
      \State $length = 0$
      \Function{VisibleLines}{$L, n$}
        \If{$n \le 2$}
          \State \Return $L$
        \EndIf
        \State $mid \gets \lfloor n/2 \rfloor$
        \State $left = \textsc{VisibleLines}(L[0:mid-1], mid)$
        \State $right = \textsc{VisibleLines}(L[mid:n-1], n - mid)$
        \State $i \gets mid - 1$
        \State $j \gets 0$
        \While{$i \ge 0$ or $j < n - mid$}
          \Comment Merge
          \If{$i = 0$}
            \State $V[length] \gets L[j]$
            \State $length \gets length + 1$
            \State $j \gets j + 1$
          \ElsIf{$j = n - mid$}
            \State $V[length] \gets L[i]$
            \State $length \gets length + 1$
            \State $i \gets i - 1$
          \EndIf
          \If{$\frac{left[i-1][1] - left[i][1]}{left[i][0] - left[i-1][0]} < \frac{right[j+1][1] - left[j][1]}{left[i][0] - right[j][0]}$}
            \State $V[length] \gets L[i]$
            \State $length \gets length + 1$
            \State $i \gets i - 1$
          \EndIf
          \If{$\frac{left[i][1] - right[j][1]}{right[j][0] - left[i][0]} < \frac{right[j+1][1] - right[j][1]}{right[j][0] - right[j+1][0]}$}
            \State $V[length] \gets L[j]$
            \State $length \gets length + 1$
            \State $j \gets j + 1$
          \EndIf
        \EndWhile
        \State \Return $V$
      \EndFunction
      \State \Return \textsc{VisibleLines}($L', n'$)
    \end{algorithmic}
  \end{algorithm}

\section{Problem 5}
  Let's model this data structure using an array.
  The children of the node at index $i$ will be at indices $2i + 1$ and $2i + 2$.
  Let \textsc{MaximumSum}($T, n, i$) be the required algorithm, where $T$ is the array storing the triangle, $n$ is the number of numbers in the triangle (or length of $T$), and $i$ is used to denote the index of the current root (will be used later).
  It is easy to see the following relation:
  \[
    \textsc{MaximumSum}(T, n, i) = T[i] + \max\{
      \underbrace{\textsc{MaximumSum}(T, n, 2i+1)}_{\text{path through } child1},
      \underbrace{\textsc{MaximumSum}(T, n, 2i+2)}_{\text{path through } child2}
    \}
  \]
  The required algorithm is:
  \begin{algorithm}
    \caption{Maximum Path Sum in Triangle}
    \begin{algorithmic}[1]
      \State \textbf{Input:} Triangle as an array $T$ of size $n$
      \Comment $T[0]$ is the root of the triangle
      \State \textbf{Output:} Maximum path sum from top to bottom of the triangle
      \State Initialize an empty array $M$ of size $n$ with $M[i] = \texttt{None} ~ \forall i \in \{0, \ldots, n-1\}$
      \Function{MaximumSum}{$T, n, i$}
        % \If{$n = 1$}
        %   \State \Return $T[0]$
        % \EndIf
        \If{$M[i] = \texttt{None}$}
          \If{$2i + 1 \ge n$}
            \State $M[i] \gets T[i]$
          \ElsIf{$2i + 2 \ge n$}
            \State $M[i] \gets T[i] + \textsc{MaximumSum}(T, n, 2i+1)$
          \Else
            \State $M[i] \gets T[i] + \max\{\textsc{MaximumSum}(T, n, 2i+1), \textsc{MaximumSum}(T, n, 2i+2)\}$
          \EndIf
        \EndIf
        \State \Return $M[i]$
      \EndFunction
      \State \Return \textsc{MaximumSum}($T, n, 0$)
    \end{algorithmic}
  \end{algorithm}

\end{document}
